\scanpage{156}
\margin{155}
\begin{center}- 112 -\end{center}\hfill IX

Revenant à la définition de (9.4.2) lorsque $S$ est strictement local, on note que la relation $\operatorname{Ext}^2(Q,P)=0$ implique la possibilité de définir un élément canonique (9.3.13), où on prend pour $Y'$ l'extension panachée $X_\eta$ donnée de $E_\eta$ par $F_\eta$. La formule (9.4.5) permet d'identifier cet élément $c(X_\eta)$ à un élément de $\check Q_{\unsure{o|0}}\otimes\check Q_{\unsure{o|0}}^*$, ce qui précise la définition de (9.4.2). On trouve de plus, en vertu des considérations générales de 9.3, que l'élément en question s'interprète comme l'obstruction au prolongement de $X_\eta$ en une extension panachée $X$ de $E$ par $F$.

9.4.6. La définition de (9.4.1) dans le cas $S$ hensélien se ramène maintenant immédiatement au cas strictement local, par descente galoisienne à partir du hensélisé strict de $S$.

\uline{Proposition} 9.4.7. \uline{Les données} $P,Q,R,E,F$ \uline{sont comme ci-dessus}, \uline{sauf qu'il est inutile de supposer le trait} $S$ \uline{hensélien}. \uline{On considère le foncteur}
\figure{\begin{tikzpicture}\draw (0,.3)--(.12,0);\end{tikzpicture}}
\[
 (9.4.7.1)\qquad X\longmapsto X_\eta
\]
\uline{de la catégorie des extensions panachées de} $E$ \uline{par} $F$ \uline{dans la catégorie des extensions panachées de} $E_\eta$ \uline{par} $F_\eta$. \uline{Pour toute extension panachée} $X_\eta$ \uline{de} $E_\eta$ \uline{par} $F_\eta$, \uline{on considère l'accouplement}
\[
 (9.4.7.2)\qquad c_{\unsure{o|0}}(X_\eta):Q_{\unsure{o|0}}\otimes Q_{\unsure{o|0}}^*\longrightarrow\Lambda_{\unsure{s|S}}
\]
\uline{induit par l'accouplement du type} (9.4.1) \uline{défini sur le} hensélisé $S^h$ \uline{de} $S$.

\uline{Ceci posé}, \uline{le foncteur} (9.4.7.1) \uline{est} pleinement fidèle, \uline{et pour une extension panachée} $X_\eta$ \uline{de} $E_\eta$ \uline{par} $F_\eta$, $X_\eta$ \uline{appartient à l'image essentielle de ce foncteur si et seulement si l'accouplement} (9.4.7.2) \uline{est nul}.
\margin{\diagram{\begin{tikzcd}\unsure{\otimes^i k}\arrow[r]&\unsure{\wedge}\end{tikzcd}}}
\margin{430}
\margin{430}
